
\section{Contexto de Plasmas Relativistas y Astrofísicos}

\primerparrafo El estudio de los fluidos ionizados en presencia de campos magnéticos, conocidos como plasmas, constituye un pilar fundamental en la física moderna, dado que más del 90\% de la materia bariónica en el Universo se encuentra en este estado \parencite{goedbloed-2004A}. La dinámica de estos sistemas, particularmente en entornos de alta energía, requiere la aplicación de marcos teóricos que incorporen los efectos de la relatividad especial y general \parencite{rezzolla-2013}. La Magnetohidrodinámica (MHD) ofrece una descripción macroscópica de los plasmas, la cual es esencial para comprender fenómenos astrofísicos asociados con plasmas relativistas fuertemente magnetizados, tales como Núcleos Galácticos Activos (AGN), chorros relativistas, vientos de púlsares, Brotes de Rayos Gamma (GRBs) y magnetares \parencite{goedbloed-2004A, rezzolla-2013}.\sebnota{agregar lo que haga falta de:enfatizar más en las aplicaciones en la astrofísica, por ejemplo: La inestabilidad de Kelvin–Helmholtz (KHI) aparece de manera natural en regiones donde
dos flujos de plasma se desplazan con velocidades diferentes. En astrofísica, este tipo de
inestabilidad es relevante en jets relativistas, vientos magnetizados, eyecciones compactas y
zonas de interacción entre plasmas, donde puede favorecer la mezcla de materia, la formación
de vórtices, la generación de estructuras secundarias y la reorganización del campo magnético.     (esto hecho con ia, pero pueden humanizarlo o reescribirlo)}

\subsection{Relevancia de la RRMHD en Chorros Astrofísicos y AGN}

\primerparrafo La aproximación más común para modelar estos sistemas es la Magnetohidrodinámica Relativista Ideal (RMHD), la cual asume una conductividad eléctrica infinita ($\sigma \to \infty$) \parencite{palenzuela-2009}. Esta aproximación es adecuada para describir las propiedades y la dinámica global de muchos fenómenos astrofísicos relativistas. Sin embargo, existen contextos donde la conductividad finita (resistividad) juega un papel crucial, haciendo necesaria la inclusión de la Magnetohidrodinámica Relativista Resistiva (RRMHD) \parencite{palenzuela-2009, miranda-aranguren-2018}. La RRMHD se vuelve imprescindible en regiones de alta disipación, como las láminas de corriente (\textit{current sheets})\footnote{Una lámina de corriente es una región delgada donde el campo magnético invierte o cambia bruscamente su dirección, concentrando una intensa densidad de corriente $\mathbf{J}=\nabla\times\mathbf{B}$; constituyen los sitios privilegiados de la reconexión magnética y de la disipación óhmica.} y las zonas turbulentas, que son energéticamente importantes y están prohibidas dentro del límite ideal debido a la conservación del flujo magnético \parencite{palenzuela-2009, miranda-aranguren-2018}. La relevancia de la RRMHD en chorros y AGN radica en la necesidad de modelar fenómenos de reconexión magnética y la disipación de energía \parencite{palenzuela-2009, miranda-aranguren-2018}.

Desde el punto de vista numérico, la inclusión de la resistividad representa un desafío considerable, ya que el sistema de ecuaciones se vuelve rígido (\textit{stiff})\footnote{Un sistema es \emph{rígido} (\textit{stiff}) cuando coexisten escalas temporales muy dispares. A alta conductividad, el tiempo de relajación resistivo $\tau\sim1/\sigma$ es órdenes de magnitud menor que la escala dinámica de advección, lo que obligaría a un integrador explícito a usar pasos de tiempo prohibitivamente pequeños.} en el límite cercano al RMHD ideal (alta conductividad) \parencite{mizuno-2013, palenzuela-2009, miranda2014}. Esto se debe a que los términos fuente para la evolución del campo eléctrico se vuelven rígidos \parencite{mizuno-2013}. Para tratar eficazmente estos sistemas se requieren métodos numéricos especializados, como los esquemas Implícito-Explícito (IMEX) Runge-Kutta (IMEX-RK) \parencite{palenzuela-2009, miranda2014}.

\subsection{La Inestabilidad de Kelvin-Helmholtz (KHI)}

\primerparrafo Dentro de la dinámica de los fluidos magnetizados, la Inestabilidad de Kelvin-Helmholtz (KHI) emerge como un mecanismo fundamental, desencadenado por una velocidad de cizallamiento (velocity shear) entre dos capas de fluido o plasma en contacto \parencite{gomes-2017, tian-2016}. La KHI, cuyo estudio clásico se remonta a los teoremas de Kelvin y Helmholtz \parencite{acheson-1990}, es crucial para el transporte de momento y energía y la mezcla de fluidos \parencite{gomes-2017, tian-2016}. En el contexto astrofísico, la KHI es un mecanismo capaz de amplificar campos magnéticos al transformar la energía cinética del flujo de cizallamiento en energía magnética \parencite{pimentel-2019}. El análisis lineal de la KHI en el régimen relativista magnetizado indica que se necesita una baja razón de la velocidad de Alfvén\footnote{La velocidad de Alfvén $V_A=B/\sqrt{\rho h+B^2}$ (en unidades con $c=1$) es la rapidez de propagación de las perturbaciones magnéticas a lo largo de las líneas de campo; cuantifica la rigidez magnética del plasma.} a la velocidad del sonido del fluido ($V_A/c_s$, donde $c_s$ es la velocidad del sonido térmica del gas, no la magnetosónica) para activar los modos inestables \parencite{osmanov-2008}.

El estudio de la KHI bajo el marco de la RRMHD es esencial para lograr una modelización precisa y realista de la dinámica de los plasmas astrofísicos. Además, la KHI, debido a su complejidad, es un problema de prueba bien conocido para la verificación de los códigos numéricos de MHD y RRMHD, incluyendo simulaciones de muy alto orden \parencite{gomes-2017}.

\section{Transición de RMHD a RRMHD}

\primerparrafo Como se mencionó anteriormente, existen situaciones donde la aproximación del RMHD (ideal) deja de tener validez, ya que hay fenómenos donde la hipótesis de conductividad infinita $(\sigma \rightarrow \infty)$ entra en conflicto con las observaciones astronómicas y con la consistencia teórica de los modelos basados en esta hipótesis. Una evidencia de que este modelo es limitado para algunas situaciones es el problema de la reconexión magnética y la topología del campo magnético. El RMHD ideal impone la condición de ``congelamiento'' del fluido, la cual consiste en que las líneas de campo magnético están ligadas al fluido y no se permiten cambios de topología del campo. Cuando se presentan fenómenos de reconexión en simulaciones de RMHD, es a causa de errores de truncamiento del esquema, lo que se llama ``resistividad numérica'' y corresponde a un proceso que no es controlado y depende exclusivamente del marco numérico \parencite{komissarov-2007, mizuno-2013}. Sin embargo, al incluir el término de resistividad $\eta$, permite que la reconexión aparezca de forma natural, predictiva y, sobre todo, física, lo cual es esencial en fenómenos como las llamaradas en AGNs, GRBs y magnetosferas de púlsares \parencite{miranda-aranguren-2018, palenzuela-2009}.

Otro fenómeno en el cual es inevitable la inclusión de un término resistivo físico es la presencia de capas de corriente (current sheets). Comparando cómo se comporta el RMHD ideal y el RRMHD en regiones con altos gradientes espaciales de campo magnético, el término disipativo asociado a la resistividad $(\eta \nabla^2 \mathbf{B})$ se vuelve dominante incluso cuando se tienen valores bajos de conductividad, por lo cual la incorporación del término resistivo es necesaria y permite que los fenómenos de aniquilación magnética\footnote{La aniquilación magnética es la cancelación de flujo magnético de polaridades opuestas que se reconectan, con la consiguiente conversión de energía magnética en energía térmica y cinética.}, conversión de energía magnética en calor y efectos dinámicos sobre las partículas sean posibles.


\section{Objetivos y Estructura de la Monografía}

\primerparrafo El trabajo en cuestión consta inicialmente de un recorrido conceptual y teórico acerca del paradigma sobre el cual está estructurada la Magnetohidrodinámica Relativista Resistiva (RRMHD), abordando tanto la física de los escenarios que representa como el fundamento teórico y matemático que respalda su formulación. Posteriormente, se trata en detalle todo lo relacionado con la inestabilidad de Kelvin–Helmholtz, desde su solución en el régimen lineal dentro de la mecánica de fluidos clásica hasta la forma en que este problema se escala y formula en el marco de la RRMHD, explorando las soluciones analíticas aproximadas disponibles en los límites accesibles (hidrodinámico clásico y RMHD).

Todo esto se realiza con el fin de estudiar la relación entre la dinámica temporal de la inestabilidad y la resistividad utilizada, mediante el uso del código \textit{Cueva}, el cual resuelve numéricamente las ecuaciones de la RRMHD en la formulación de \textcite{komissarov-2007}, implementada con el solucionador de Riemann HLLC resistivo desarrollado por \textcite{miranda-aranguren-2018}. En este contexto, se cuantifica cómo la resistividad modifica la tasa de crecimiento de la inestabilidad, se analiza su efecto en la formación de estructuras secundarias, como los vórtices, y se evalúa cómo la difusión magnética causada por la resistividad afecta el crecimiento de la inestabilidad y la evolución de las variables electromagnéticas globales (corriente, energía magnética y disipación) en toda la malla de simulación.

Finalmente, se evalúa la forma en la que las variaciones espaciales, tanto en dirección como en magnitud, del campo magnético externo modifican la evolución temporal de dicha inestabilidad en la interfaz de corte.\revnota{este objetivo (orientación/intensidad de B) NO se desarrolla en el cuerpo de esta versión; es la campaña pendiente (ver Conclusiones, objetivo 4). Alinear la redacción a ``se sientan las bases para evaluar\dots'' o presentarlo como trabajo futuro.} El objetivo de comprender estas relaciones es poder explicar, con mayor detalle físico, cómo los fenómenos astrofísicos de plasmas relativistas en los cuales este tipo de inestabilidades es ubicuo se ven afectados por la influencia de la resistividad.